\section{Implementación del software}
\subsection{C. Fase 3- Diseño y desarrollo del sistema}
Transformar los criterios teóricos y metodológicos en un sistema funcional de gestión de eventos, aplicando la arquitectura MVC, la POO y la metodología ágil Scrum como guías de construcción y organización del proyecto.

\subsubsection{Procedimiento detallado}
\paragraph{1. Diseño de la arquitectura (MVC):}
• Se definieron las capas principales del sistema: Modelo (gestión de datos y lógica de negocio), Vista (interfaces gráficas de usuario) y Controlador (coordinación de entradas/salidas y flujo de información).

• Esta separación permitió mejorar la mantenibilidad y la escalabilidad del sistema, ya que cada capa pudo evolucionar de manera independiente.

\paragraph{2. Programación Orientada a Objetos (POO):}
• Se estructuró el código bajo principios de herencia, encapsulamiento y polimorfismo, lo cual garantizó un desarrollo más ordenado y reutilizable.

• El uso de POO facilitó la construcción de módulos comunes para diferentes funcionalidades, como gestión de usuarios, administración de eventos y reportes.

\paragraph{3. Metodología Ágil Scrum:}
• Se realizaron reuniones periódicas de seguimiento y retroalimentación, permitiendo ajustar prioridades y detectar errores de manera temprana.

\paragraph{4. Herramientas utilizadas:}
• Backend: Java con Spring Boot, seleccionado por su robustez, seguridad y escalabilidad.

• Frontend web: React.js, elegido por su flexibilidad en el diseño de interfaces dinámicas.

• Aplicación móvil: React Native, lo que permitió extender las funcionalidades al entorno móvil sin duplicar esfuerzos de programación.

• Base de datos: PostgreSQL, empleada para gestionar la información de usuarios, eventos y reservas.

• Control de versiones: GitHub, para trabajo colaborativo y control de cambios en el código.